\input{../000_header_year.txt}

%%%%%%%%%%%%%%%%%%%%%%%%%%%%%%%%%%%%%%%%%%%%%
%　遷移行列を求める問題
%　MATLAB mファイルでA行列を指定して e^At を求める
%
%　mファイルの使い方
%　　　まず A 行列を定義する
%　　　固有ベクトル整数とするために p1, p2を試行錯誤で修正
%
%%%%%%%%%%%%%%%%%%%%%%%%%%%%%%%%%%%%%%%%%%%%%


\begin{document}


%++++++++++++++++++++++++++++++++++++++++++++++++++++++++++++++++++++++++++++++++
%++++++++++++++++++++++++++++++++++++++++++++++++++++++++++++++++++++++++++++++++
%++++++++++++++++++++++++++++++++++++++++++++++++++++++++++++++++++++++++++++++++
\exheader
%++++++++++++++++++++++++++++++++++++++++++++++++++++++++++++++++++++++++++++++++
%++++++++++++++++++++++++++++++++++++++++++++++++++++++++++++++++++++++++++++++++
%++++++++++++++++++++++++++++++++++++++++++++++++++++++++++++++++++++++++++++++++

%================================
%	matlab060.mを利用
%	パラメータの定義
%		A行列
%		p1,p2それぞれ第１，第２固有ベクトルをわかりやすい値にするために乗ずる係数
%
%	chrA:A行列
%	chrM0: lambda I -A
%	chrlambda1, charlambda2: 固有値
%	chreqn: 特性多項式
%	chrM1: 第１固有値を代入した lambda I - A
%	chrV1: 第１固有ベクトル（第１要素を１とした後 p1を乗じたもの）
%	chrM2: 第２固有値を代入した lambda I - A
%	chrV2: 第２固有ベクトル（第１要素を１とした後 p2を乗じたもの）
%	chrT: 対角化のための行列T
%	chrTI: Tの逆行列
%	chreberA: 対角化した後の遷移行列
%	chreA: 遷移行列
%


\input{065_def_file.txt}

%================================

\noindent
{\bf 遷移行列（演習）} \nopagebreak

次の行列$A$の遷移行列$e^{At}$を求めよ．
\[
	A= \chrA
\]
hint: 行列$A$を$\Lambda = T^{-1}AT$と対角化し，$e^{\Lambda t}$が解析的に求まることを利用する．


\vfill

\begin{tabular}{ll}
	\hspace*{8cm} &
	\begin{tabular}{|l|}
		\hline 
		解答：遷移行列\\
		\hspace*{8cm} \\[2cm]
		\hline
	\end{tabular}
\end{tabular}


\newpage

%---------------------------------------------------------------------------


\noindent
{\bf 遷移行列（演習）} \nopagebreak

次の行列$A$の遷移行列$e^{At}$を求めよ．
\[
	A= \chrA
\]
hint: 行列$A$を$\Lambda = T^{-1}AT$と対角化し，$e^{\Lambda t}$が解析的に求まることを利用する．\\

$A$の固有値は，
$\det(\lambda I - A) = 0$を満たす$\lambda$，すなわち
\begin{align*}
 \det(\lambda I - A)
 & = 
 \det \chrMo
 =
 (\chrMoa)(\chrMod)-(\chrMob)(\chrMoc)
\\
 &=
 \chreqn
 =
 \chreqno
\end{align*}
より，$\lambda = \chrlambdaa または \chrlambdab$である．

$\lambda = \chrlambdaa$のとき，
\[
 \begin{array}{ccc}
 (\lambda I - A) v = 0 
 &\Leftrightarrow&
 \chrMa v = 0 
 \end{array}
\]
であるから，$\lambda = \chrlambdaa$に対する固有ベクトルは
$v_1 = \chrVa$．

同様に
$\lambda = \chrlambdab$のとき，
\[
 \begin{array}{ccc}
 (\lambda I - A) v = 0
 &\Leftrightarrow&
 \chrMb v = 0 
 \end{array}
\]
であるから，$\lambda = \chrlambdab$に対する固有ベクトル.\
$v_2 = \chrVb$.

$A$を対角化する行列は
\[
 T=(v_1, v_2)
  =
 \chrT
 , \ \ 
 T^{-1}= \chrTI
\]
となる．$A$を対角化すれば
\[
 \Lambda = T^{-1}AT 
 = \chrTI \chrA \chrT
 = \chrbarA
\]
となる．よって
\begin{align*}
 e^{At} &= T e^{T^{-1}AT t} T^{-1}
 = T e^{\Lambda t} T^{-1}
\\
 &= \chrT \chrebarA \chrTI
 = \chreA
\end{align*}



\end{document}

